% Includable supplement. Independent of the proposed arithmetic construction.
\section{A restriction on the degree in two witnesses}
\label{sec:p2lb-conics}

The total degree can conceal an important distinction: a polynomial may
have arbitrarily high degree in the input, while remaining quadratic in
its witnesses.  With only two witnesses, such a polynomial cannot define
the ring of integers.  We prove a simultaneous version which also rules
out a finite disjunction of such equations.

\begin{theorem}\label{p2lb:quadratic-witnesses}
Let $K$ be a number field and let
$F_1,\ldots,F_s\in K[T,Y,Z]$ have total degree at most two in the
pair $(Y,Z)$, with no restriction on their degree in $T$.  Put
\[
 S=\{t\in K:\text{for some }i\text{ and }y,z\in K,
                         \ F_i(t,y,z)=0\}.
\]
Exactly one of the following alternatives holds:
\begin{enumerate}
\item at least one equation has a solution
$(Y,Z)\in K(T)^2$, and $S$ is cofinite in $K$;
\item no equation has such a solution, and $K\setminus S$ contains
infinitely many elements of every coset $a+\mathfrak a$, where
$a\in\mathcal O_K$ and $\mathfrak a$ is a nonzero ideal of
$\mathcal O_K$.
\end{enumerate}
In particular, no single polynomial quadratic in two witnesses, or
finite disjunction of such polynomials, defines $\mathcal O_K$ in $K$.
\end{theorem}

The cofinite assertion in the first alternative follows by specializing
the rational functions in a generic solution away from their finitely
many poles.  The substantive issue is to find integral specializations
which simultaneously preserve the obstructions in the second alternative.
We use Faddeev's residue sequence, Chebotarev's theorem and integral
Hilbert irreducibility.  The argument is a consequence of these classical
tools; no claim of historical priority is made.

\subsection{Quaternion obstructions at integral specializations}

For a field $F$ of characteristic different from two, $(a,b)$ denotes the
Brauer class of the quaternion algebra generated by $i,j$ with
$i^2=a$, $j^2=b$ and $ij=-ji$.  A smooth conic is split, or equivalently
has an $F$-point, precisely when its associated quaternion class is zero.
For this correspondence and the residue sequence below see
\cite[Theorem~5.1 and Section~6.1]{p2lbBergEtAl2024}.

We recall explicitly the residue convention.  At the closed point of
$\mathbb A^1_K$ given by a monic irreducible $f(T)$, write
\[
 a=f^r u,\qquad b=f^s v,
\]
where $u,v$ are units in $K[T]_{(f)}$.  If $\theta$ is a root of $f$
and $L=K(\theta)$, the residue of $(a,b)$ is the square class of
\[
             (-1)^{rs}u(\theta)^s/v(\theta)^r
             \quad\text{in }L^\times/L^{\times2}.
\]
Faddeev's theorem implies that a class in $\operatorname{Br}K(T)$
with all finite residues zero comes from $\operatorname{Br}K$.
Indeed the sum-of-corestrictions relation first makes the residue at
infinity zero, since its residue field is $K$; exactness then gives a
constant class.  In characteristic zero the sequence applies to the
whole Brauer group because $\operatorname{Br}\overline K(T)=0$.

We shall also use the following elementary local formula.  If $F$ is a
nonarchimedean local field with odd residue field $\kappa$, if $\varpi$
is a uniformizer, and if $u,v$ are units, the quaternion
$(\varpi^r u,\varpi^s v)$ is split exactly when
\[
                    (-1)^{rs}\overline u^{\,s}/\overline v^{\,r}
                    \quad\text{is a square in }\kappa^	imes.
\]
One obtains this by bilinearity of the Hilbert symbol: two unit
parameters give a split algebra, $(\varpi,u)$ is detected by the square
class of $\overline u$, and $(\varpi,\varpi)=(\varpi,-1)$.
These facts also follow from the norm criterion for the unramified
quadratic extension; see \cite[Proposition~5.5]{p2lbVoight2019}.

\begin{lemma}\label{p2lb:brauer-progression}
Let $\alpha\in\operatorname{Br}K(T)$ be a nonzero quaternion class.
If $\alpha$ is constant, every specialization outside a finite set
is nonzero.  If $\alpha$ is not constant, then outside any prescribed
finite set of primes of $K$ there is a prime $\mathfrak p$ and a
residue class $c\bmod\mathfrak p^2$ such that, outside a finite set
of inputs, every $t\in\mathcal O_K$ in that residue class has
\[
                   \alpha(t)\ne0\quad\text{in }
                   \operatorname{Br}K_{\mathfrak p}.
\]
\end{lemma}
\begin{proof}
The assertion for a constant class follows from specialization on an
open set where its equality with the chosen quaternion presentation is
defined.  Otherwise choose a nonzero finite residue using Faddeev's
theorem.  With the notation just introduced, let
\[
 d=(-1)^{rs}u(\theta)^s/v(\theta)^r\in L^\times
\]
be its nonsquare representative.  Choose a finite Galois extension
$M/K$ containing $L(\sqrt d)$, and an element
$g\in\operatorname{Gal}(M/L)$ which sends $\sqrt d$ to $-\sqrt d$.
Such an element exists because $d$ is not a square in $L$.

Chebotarev's theorem gives infinitely many primes of $K$ with Frobenius
conjugate to $g$; see \cite[Section~2.1, Theorem~1]{p2lbSerre1981}.
Choose one outside the prescribed finite set and all further primes of
bad reduction for the displayed algebraic data, including primes above
two.  Choose a prime of $M$ whose Frobenius is exactly $g$, and let
$\mathfrak q$ be its restriction to $L$.  Since $g$ fixes $L$, this
prime of $L$ has residue degree one over $\mathfrak p$.  Since $g$
acts nontrivially on $\sqrt d$, the reduction of $d$ at
$\mathfrak q$ is a nonsquare.  The excluded bad primes ensure that
it is a nonzero unit.

The corresponding completion of $L$ is $K_{\mathfrak p}$, so
$f$ has a simple root $\theta_{\mathfrak p}$ in
$\mathcal O_{K_{\mathfrak p}}$.  Choose a residue class
$c\bmod\mathfrak p^2$ for which
\[
 v_{\mathfrak p}(t-\theta_{\mathfrak p})=1
 \qquad(t\equiv c\bmod\mathfrak p^2).
\]
There is such a class because the reduction map from $\mathcal O_K$
onto $\mathcal O_{K_{\mathfrak p}}/\mathfrak p^2$ is surjective.
For all these inputs, $v_{\mathfrak p}(f(t))=1$, and $u(t),v(t)$
are units whose reductions equal those at $\theta$ and
$\mathfrak q$.  The local formula therefore makes the specialized
quaternion nonsplit, since its residue square class is that of $d$.
The finitely many inputs where a rational presentation is undefined may
be discarded.  This proves the assertion.
\end{proof}

\begin{corollary}\label{p2lb:simultaneous-conics}
Let $C_1,\ldots,C_m$ be smooth conics over $K(T)$ with no
$K(T)$-point.  Given any coset $a+\mathfrak a$ as in
Theorem~\ref{p2lb:quadratic-witnesses}, it contains a smaller nonempty
coset of a nonzero ideal such that, outside finitely many inputs, none
of the specialized conics has a $K$-point.
\end{corollary}
\begin{proof}
Constant nonzero quaternion classes already exclude every good
specialization.  For each remaining conic, use the preceding lemma
with a fresh prime, distinct from all previously chosen primes and
coprime to $\mathfrak a$.  The Chinese remainder theorem simultaneously
imposes its bad residue class modulo the square of that prime and the
original coset.  Every resulting good specialization is locally insoluble
at the selected prime for each corresponding conic.
\end{proof}

\subsection{Degenerate conics and the simultaneous theorem}

\begin{proof}[Proof of Theorem~\ref{p2lb:quadratic-witnesses}]
Write $F=K(T)$.  If one equation has a point in $F^2$, its rational
functions specialize to a point at every input outside a finite set.
This gives the first alternative.  Suppose henceforth that none does.
We classify each equation over $F$, and discard finitely many inputs
so that all rational coordinate changes, ranks and nonzero coefficients
in this classification specialize correctly.

An equation of witness degree zero is a nonzero element of $F$, and
has no point at any good input.  Witness degree one always gives an
affine $F$-point, so cannot occur in the case under consideration.
For witness degree two, homogenize in $Y,Z$ with a third coordinate
$W$, giving a ternary quadratic form and its projective conic $C$.
If the form has rank three and its quaternion class is zero, then
$C\simeq\mathbb P^1_F$.  The affine open $W\ne0$ is nonempty
and hence has an $F$-point, a contradiction.  Thus every rank-three
case has a nonzero quaternion class, and
Corollary~\ref{p2lb:simultaneous-conics} simultaneously excludes all
of them on a smaller coset of the given integral progression.

A rank-one ternary form is a scalar times the square of an
$F$-linear form.  That line cannot be the line at infinity, since the
original polynomial has witness degree two.  It therefore has an
affine $F$-point, again a contradiction.  A split rank-two form is the
product of two distinct $F$-linear forms.  At least one line is not
the line at infinity, and supplies the same contradiction.

The remaining possibility is a nonsplit rank-two form.  After a linear
change over $F$ it is a nonzero scalar times
\[
                         U^2-d(T)V^2,
                         \qquad d(T)\notin F^{\times2}.
\]
Its only projective $F$-point is its radical point, defined by
$U=V=0$.  If that point were affine, it would give an $F$-point of
the original equation.  Therefore its original $W$ coordinate is
zero.  At every good specialization for which $d(t)$ remains a
nonsquare, the same reasoning leaves only this radical point, still
at infinity, and the specialized affine equation has no solution.

There are only finitely many nonsquare functions $d(T)$ to treat.
The smaller coset already obtained contains a progression
$t=c+bz$ with $c\in\mathcal O_K$, $0\ne b\in\mathcal O_K$
and $z\in\mathcal O_K$: choose $b$ in its defining nonzero ideal.
The polynomials
\[
                         Y^2-d(c+bZ)
\]
remain irreducible over $K(Z)$, because $T\mapsto c+bZ$ is an
isomorphism of rational function fields.  The integral Hilbert
irreducibility theorem, simultaneously for this finite list, gives
infinitely many integral $z$ for which all these specializations
remain irreducible and all previously excluded finite conditions are
avoided.  See the integral form recalled in the original lower-bound
section and \cite[Introduction]{lbMuller2001}.  If the list is empty,
any infinitely many good inputs in the coset suffice.

The resulting inputs lie in the original coset $a+\mathfrak a$,
and none admits a solution of any equation.  This proves the second
alternative and the theorem.
\end{proof}

\begin{corollary}\label{p2lb:two-witness-degree-three}
If $P(T,Y,Z)$ defines $\mathbb Z$ in $\mathbb Q$, then its total
degree in the two witnesses $Y,Z$ is at least three.  More generally,
after factoring over $\mathbb Q(T)$, at least one factor has total
witness degree at least three.  This remains true after every
invertible affine change of the witness coordinates over
$\mathbb Q(T)$, provided the transformed equation and inverse change
are defined at all but finitely many input values.
\end{corollary}
\begin{proof}
The first assertion is the theorem with one equation.  If all factors
had witness degree at most two, clear their rational-function
denominators.  Outside finitely many inputs, the original zero condition
is their finite disjunction, contradicting the theorem.  An invertible
affine change over $\mathbb Q(T)$ preserves the rational solution sets
at every good specialization, so the same argument applies after it.
\end{proof}

\begin{corollary}\label{p2lb:conic-integral-superset}
Let $S\subset K$ contain $\mathcal O_K$, and suppose that $K\setminus S$
is infinite.  Then $S$ is not a finite union of projections of equations
in two witnesses which are jointly quadratic in those witnesses, even if
the degree in the input is unrestricted.
\end{corollary}
\begin{proof}
The second alternative of Theorem~\ref{p2lb:quadratic-witnesses} would
omit infinitely many elements of $\mathcal O_K$, while the first would
make $K\setminus S$ finite.  Both contradict the stated assumptions.
\end{proof}

For example, take
\[
 S=\mathbb Z_{(2)}\cap\mathbb Z_{(3)}
   =\{a/b\in\mathbb Q:\gcd(a,b)=1,\ b>0,\ 2\nmid b,\ 3\nmid b\}.
\]
This set contains $\mathbb Z$ and has infinite complement, for instance
the distinct elements $2^{-n}$ for $n\ge1$.  Thus a quadratic
three-witness definition of this semilocal ring cannot be replaced by a
quadratic two-witness definition, or a finite union of such definitions.
The obstruction concerns joint witness degree; it allows arbitrary input
degree and does not assert a lower bound for equations of higher witness
degree.

This is an unconditional restriction on a possible two-witness
definition at every input degree.  It does not exclude general quartic
surfaces, whose generic fibers over the input can have positive genus.
It also does not extend by this proof to quadratic equations in three
witnesses: their fibers are quadric surfaces, and are not controlled by
a single quaternion obstruction of the kind used above.
